% Analogue switch (jfet-switch): the set-up a person would wire, semantically the deck's.
%   VS s 0 1          the signal source, on the source terminal          -> source V_S (set)
%   JW l g s DUT      the switch, drain to the load                      -> njfet
%   RL l 0 1000       the load                                           -> resistor R_L
%   RG in g 1000, VIN in 0 PWL(…)   the gate drive through 1 kΩ          -> resistor R_G + source V_in (set)
%   V(l) read                       the load voltage                     -> voltmeter v_load (measured)
%   The charge-injection deck replaces RL by a 1 nF hold capacitor and puts the signal at 0 V.
\documentclass[tikz, border=10pt]{standalone}
\usepackage{bench}
\begin{document}
\begin{tikzpicture}
\benchtitle{0,7.0}{Analogue switch}
\def\gnd{1.0}
% the transistor
\node[njfet] (Q) at (4,3) {};
\node[anchor=west, font=\small, align=left] at ($(Q.D)+(0.15,-0.55)$) {DUT};
% the source terminal on the signal source
\draw (Q.S) -- ++(0,-0.3) coordinate (sj);
\draw[set] (sj) to[V, l=$V_S$] (sj |- 0,\gnd) node[ground, black]{};
\nodename{anchor=east}{$(Q.S)+(-0.1,-0.2)$}{s}
% the drain to the load node, the load down to ground, the voltmeter beside it
\draw (Q.D) -- ++(0,0.8) coordinate (l) -- ++(2.2,0) coordinate (rl);
\draw (rl) to[R, l_=$R_L$] (rl |- 0,\gnd) node[ground]{};
\nodename{anchor=south}{$(l)+(1.0,0.05)$}{l}
% the voltmeter on rl
\coordinate (vm) at ($(rl)+(1.8,0)$);
\draw (rl) -- (vm);
\voltmeter{vm |- 0,\gnd}{vm}{$v_{load}$}
\draw (vm |- 0,\gnd) node[ground]{};
% the gate resistor and the pulse source
\coordinate (g) at (Q.G);
\draw (g) -- ++(-0.3,0) coordinate (gj) to[R, l_=$R_G$] ++(-1.8,0) coordinate (vin);
\draw[set] (vin) to[V, l_=$V_{in}$] (vin |- 0,\gnd) node[ground, black]{};
\nodename{anchor=south}{$(g)+(-0.15,0.05)$}{g}
\nodename{anchor=south east}{$(vin)+(-0.05,0.05)$}{in}
% the definitions
\definitions{0,-0.1}{%
  $V_S = 1$ V;\quad $R_L = 1$ k\ohm{} and 10 k\ohm, one deck each;\quad $R_G = 1$ k\ohm\\[2pt]
  $V_{in}$: $V_{off}$, then 0 V from 100 ns to 2.11 \textmu s, then $V_{off}$;\quad 10 ns edges;\quad $V_{off} = -10$ V, or twice the model's own $V_{GS(off)}$\\[2pt]
  read $v_{load}(t)$: $t_{on}$ from the gate's rising edge to the 90\,\% crossing, $t_{off}$ from the falling edge to the 10\,\% crossing\\[2pt]
  charge injection: $V_S = 0$, $R_L$ replaced by a 1 nF hold capacitor, the gate stepped from 0 V to $V_{off}$ once;\quad $Q = C$ times the step of $v_{load}$%ENDDEFS
}
\end{tikzpicture}
\end{document}
